\section{Part III: Discovering Problems That Are Worth Solving}

\subsection{The Hierarchy of Problem Evidence}

A founder hears an engineer say, ``I struggle with evaluating AI models,'' and rushes to declare a validated startup opportunity. A verbal complaint does not establish urgency, willingness to sacrifice, or purchasing authority. We must distinguish 7 distinct rungs of problem evidence:
\begin{enumerate}[leftmargin=1.2em,itemsep=0.2ex]
  \item \textbf{Abstract opinion:} General agreement that a problem is interesting or culturally important.
  \item \textbf{Specific recent incident:} Detailed recollection of a concrete event where the problem caused friction.
  \item \textbf{Measurable consequence:} Documented hours lost, revenue delayed, or software crashes incurred.
  \item \textbf{Costly workaround:} Active expenditure of personal time, internal labor, or makeshift software scripts.
  \item \textbf{Active search or budget request:} Verifiable steps taken to evaluate commercial solutions.
  \item \textbf{Direct financial expenditure:} Historical payments made to alternative vendors or consultants.
  \item \textbf{Repeatable, full-price purchase:} Multiple independent buyers paying sustainable prices and realizing value.
\end{enumerate}

Evidence on rungs 1 and 2 is helpful for learning customer vocabulary; only rungs 4 through 7 indicate \textbf{commercially actionable demand}.

\subsection{The Anatomy of Acute Demand}

An acute problem exhibits a specific combination of structural attributes:
\begin{itemize}[leftmargin=1.2em,itemsep=0.2ex]
  \item \textbf{Severity:} Failure imposes measurable economic, operational, or status damage.
  \item \textbf{Time sensitivity:} Delay steadily compounds the damage or misses a critical external window.
  \item \textbf{Clear ownership:} A specific individual is held personally accountable for resolving the outcome.
  \item \textbf{Inadequate alternatives:} Existing free or internal tools have failed to solve the issue.
  \item \textbf{Accessible budget:} The accountable owner controls funds or has an approved procurement pathway.
\end{itemize}

Consider 3 technical learners: A hobbyist explores machine learning out of general curiosity; a developer repeatedly loses weekend evenings fixing evaluation pipelines; a tech lead must pass a strict reliability audit in 45 days or lose an enterprise client. The third case represents acute commercial demand, but that buyer may strongly prefer a 2-day architectural review service over a 6-month broad curriculum. Always ask: \textbf{acute for whom, and through which exact purchasing mechanism?}

\subsection{The Multidimensional Nature of Problem Cost}

Cost is rarely restricted to financial expenditure. A problem's true cost spans:
\begin{itemize}[leftmargin=1.2em,itemsep=0.2ex]
  \item \textbf{Direct capital:} Invoices, vendor subscriptions, and consultant retainers.
  \item \textbf{Labor time:} Highly paid engineering hours diverted from core roadmap features to rework.
  \item \textbf{Opportunity loss:} Delayed product releases and foregone commercial contracts.
  \item \textbf{Career and status risk:} Public embarrassment in executive reviews and loss of promotion velocity.
  \item \textbf{Psychological drag:} Chronic anxiety, team attrition, and developer burnout.
\end{itemize}

Ask customers for concrete bounds grounded in actual episodes: How many hours were spent last week? What feature release was delayed? Which invoice or error log illustrates the incident? Avoid inflating minor inconveniences into dramatic financial crises; not every saved hour converts into cash.

\subsection{Frequency, Severity, and Recoverability}

Problems behave differently depending upon their operational parameters:
\begin{itemize}[leftmargin=1.2em,itemsep=0.2ex]
  \item \textbf{High-frequency, low-severity:} Daily syntax lookup errors; justified by free tools, search, or low-cost subscriptions.
  \item \textbf{Low-frequency, catastrophic-severity:} Annual regulatory compliance audits or major system outages; justify high-ticket specialist interventions.
  \item \textbf{One-time inflection events:} Pivotal career transitions or foundational college curriculum updates; justify substantial one-off investments.
\end{itemize}
Irreversibility matters: an error easily undone with a unit test creates low willingness to pay; an error that permanently leaks confidential customer data creates intense willingness to pay.

\subsection{Workarounds as Proof of Costly Demand}

A workaround is unvarnished behavioral proof of demand. When customers maintain fragile, 40-tab Excel workbooks, write bespoke Python scrapers, hold ad-hoc weekend peer review sessions, or hire unvetted freelancers on personal credit cards, they are voting with their scarce time and capital.

Examine \textbf{what the workaround costs and what it protects}. A manual spreadsheet may be painfully slow, yet uniquely flexible, fully inspectable, and compliant with internal audit rules. Replacing a workaround succeeds only when the new solution preserves its hidden benefits while eliminating its primary friction.

\subsection{The Cost of Delay as a Diagnostic}

Ask: ``What happens if this problem remains completely unaddressed for 30 days, 6 months, or 1 year?'' If the honest answer is ``virtually nothing changes,'' willingness to pay is low, regardless of how enthusiastic the conversation sounds. If the answer is ``our enterprise contract is canceled,'' urgency is real.

Use naturally occurring organizational clocks: software release schedules, fiscal budget expiration dates, audit deadlines, accreditation reviews, or promotion cycles. Never invent apocalyptic scenarios that customers do not believe.

\subsection{Budget as Proof of Priority, Not Fit}

Ask which specific budget line currently finances existing workarounds: contractor fees, travel budgets, cloud compute overruns, or formal training allocations.

An existing budget proves that leadership has authorized spending for that class of problem. However, entrenched vendor relationships, multi-year contracts, and strict procurement lists can lock that money away from a newly incorporated startup. A new product category with no established budget line faces the opposite barrier: leadership may agree with the vision, but no manager possesses the authority to write the check.

\subsection{Problem Discovery Through the Study of Failure}

Success stories conceal systemic constraints through survivorship bias; failed initiatives reveal where processes break. Interview managers whose software pilots were canceled, learners who dropped out of technical bootcamps, and deans whose external partnerships collapsed.

Trace the sequence before abandonment: Was the syllabus overly theoretical? Did students lack prerequisites? Was faculty buy-in missing? Did software onboarding consume too much time? Distinguish failures caused by poor vendor execution from failures caused by fatal, unbridgeable structural mismatches.

\subsection{Badly Solved Problems vs. Unsolvable Business Models}

A customer can experience intense, genuine pain that a specific startup cannot serve at sustainable unit economics. Offering free 1-on-1 tutoring to underprivileged learners addresses an urgent societal need, but cannot finance an independent institution without philanthropic endowments or government grants. Some enterprise problems require 18-month security audits that bankrupt a seed-stage team.

The entrepreneurial objective is locating the intersection of \textbf{acute customer need, accessible budget, reachable authority, and sustainable delivery margins}. Compassion for human struggle is essential; confusing compassion with a viable commercial business model is fatal.
